\section{Background and Related Work}
\label{sec:related}

This section reviews automated test suite generation, execution-driven agentic interfaces, and the challenge of compliance bias in repository-level software maintenance.

\subsection{Large Language Models for Automated Test Generation}
Modern automated software testing techniques can be classified based on their inputs and primary operating modes. Regarding inputs, test suites have been derived from source code signatures~\cite{bareiss2022code, schafer2023empirical}, formal program specifications~\cite{docprompting}, semi-structured documentation~\cite{testexplora}, and previous program execution traces~\cite{swebench}. Black-box approaches generate test inputs directly from API specifications or documentation without requiring source code access, whereas white-box approaches analyze program structure to maximize branch and statement coverage~\cite{evosuite, randoop}.

Recent advancements in Large Language Models (LLMs) have transformed automated test generation from heuristic input fuzzing to semantic test synthesis. Tools such as ChatUnitTest~\cite{xie2023chatunitest} and CODAMOSA~\cite{lemieux2023codamosa} leverage pre-trained code models to generate unit tests, achieving high code coverage on self-contained Java methods. In repository-level contexts, DocPrompting~\cite{docprompting} and RepoBench~\cite{repobench} utilize multi-file context to complete method calls. However, all these existing techniques are primarily limited in the types of defects they can detect: they excel at identifying runtime crashes or unhandled exceptions, but struggle to detect logic flaws that go beyond mere syntax.

\subsection{Execution-Driven Agents and Dynamic Execution Signals}
Search-Based Software Testing (SBST) tools such as EvoSuite~\cite{evosuite} and Randoop~\cite{randoop} utilize genetic algorithms and random input generation to monitor program execution dynamically. While SBST tools excel at discovering edge-case crashes by executing code paths, they lack natural language understanding, producing unreadable test assertions with high false-positive rates~\cite{siddiq2023exploratory}.

To bridge the gap between high-level semantic reasoning and dynamic execution, recent work introduced agent-computer interfaces. SWE-agent~\cite{sweagent} equipped LLMs with interactive bash shell environments, enabling autonomous file navigation, tool execution, and real-time execution feedback for repository-level software engineering tasks. While SWE-agent proved effective for automated bug fixing on SWE-bench~\cite{swebench}, its application to proactive test suite generation, intent-driven assertion refinement, and compliance bias mitigation remains unexplored.

\subsection{Compliance Bias and Oracle Generation Constraints}
A common limitation of static LLM-based test generators is their vulnerability to \textit{compliance bias}~\cite{testexplora}. Because static models generate test code in a single inference pass without execution signals, they exhibit a passive tendency to conform to existing code implementations. When evaluated on defective repositories, static generators produce unit tests that syntactically conform to the buggy code, yielding a Fail-to-Pass (F2P) rate floor of only 16.06\% on the TestExplora benchmark~\cite{testexplora}. To the best of our knowledge, ours is the first approach for automated repository-level test suite generation that integrates interactive terminal execution feedback to dynamically reject compliance assumptions.
